% =========================================================================
% PLANTILLA_GENERAL_CON_EJEMPLOS.tex
% Plantilla Académica Sobria y Minimalista (Ciencias e Ingeniería)
% Configurada para el estándar de Ph.D. en Quantitative Systems Engineering
% =========================================================================
\documentclass[11pt, a4paper]{article}

% --- Paquetes Esenciales ---
\usepackage[utf8]{inputenc}
\usepackage[T1]{fontenc}
\usepackage[english, spanish, es-tabla]{babel}
\spanishdeactivate{~<>"}
\usepackage{amsmath, amssymb, amsfonts}
\usepackage{graphicx}
\usepackage{booktabs}
\usepackage{xcolor}
\usepackage{listings}
\usepackage{caption} % Personalización avanzada de leyendas (captions)

% --- Tipografía de Alta Calidad (Charter) ---
\usepackage{charter}
\usepackage{microtype}
\usepackage{metalogo}

% --- Diseño de Página y Espaciados ---
\usepackage[margin=2.5cm]{geometry}
\usepackage{setspace}
\setstretch{1.12}

% --- Paleta de Colores Sobria ---
\definecolor{primary}{rgb}{0.0, 0.35, 0.35}  % Teal apagado - Formal y sobrio
\definecolor{secondary}{rgb}{0.2, 0.2, 0.2} % Gris carbón para texto principal
\definecolor{codebg}{rgb}{0.97, 0.97, 0.95}  % Fondo suave para código
\definecolor{grayborder}{rgb}{0.85, 0.85, 0.85}

% Aplicar color secundario al cuerpo del texto
\makeatletter
\newcommand{\globalcolor}[1]{%
  \color{#1}\global\let\default@color\current@color
}
\makeatother
\AtBeginDocument{\globalcolor{secondary}}

% --- Hipervínculos Elegantes ---
\usepackage{hyperref}
\hypersetup{
    colorlinks=true,
    linkcolor=primary,
    urlcolor=primary,
    citecolor=primary,
    pdftitle={Plantilla General Académica con Ejemplos},
    pdfauthor={Jair Aguilar Peña}
}

% --- Comandos Personalizados ---
\newcommand{\vect}[1]{\mathbf{#1}}

% =========================================================================
% CONFIGURACIÓN DE JERARQUÍA DE CAPTIONS (LEYENDAS)
% =========================================================================
% Se personalizan las etiquetas y se les da un formato consistente utilizando
% el color primario de la plantilla para destacar el identificador.
\DeclareCaptionFont{teal}{\color{primary}}
\captionsetup[figure]{
    labelfont={bf,teal},
    textfont=normalfont,
    labelsep=colon,
    skip=10pt
}
\captionsetup[table]{
    labelfont={bf,teal},
    textfont=normalfont,
    labelsep=colon,
    skip=10pt
}
\captionsetup[lstlisting]{
    labelfont={bf,teal},
    textfont=normalfont,
    labelsep=colon,
    skip=10pt
}

% --- Soporte Multilingüe para Leyendas (Babel) ---
\addto\captionsspanish{%
  \renewcommand{\figurename}{Figura}%
  \renewcommand{\tablename}{Tabla}%
  \renewcommand{\lstlistingname}{Fragmento de código}%
}
\addto\captionsenglish{%
  \renewcommand{\figurename}{Figure}%
  \renewcommand{\tablename}{Table}%
  \renewcommand{\lstlistingname}{Code Snippet}%
}

% --- Guía de Expansión de Leyendas (Futuras Extensiones) ---
% Si en el futuro desea añadir nuevos elementos con sus propios prefijos 
% (por ejemplo, "Gráfico" o "Diagrama"), puede cargar el paquete `newfloat`
% o `float` en el preámbulo y declarar el nuevo tipo de la siguiente manera:
%
%   \usepackage{newfloat}
%   \DeclareFloatingEnvironment[fileext=lod, listname={Índice de Diagramas}, name=Diagrama]{diagrama}
%   \captionsetup[diagrama]{labelfont={bf,color=primary}, textfont=normalfont}
%
%   Esto le permitirá usar el entorno \begin{diagrama} ... \end{diagrama} con su propio
%   contador independiente y un caption que dirá "Diagrama X: Título".

% =========================================================================
% SOPORTE MULTILENGUAJE DE CÓDIGO (LISTINGS)
% =========================================================================
% Definición del estilo visual unificado
\lstdefinestyle{customcode}{
    backgroundcolor=\color{codebg},
    commentstyle=\color{primary!70}\itshape,
    keywordstyle=\color{primary}\bfseries,
    numberstyle=\tiny\color{gray},
    stringstyle=\color{orange!80!black},
    basicstyle=\ttfamily\small,
    breakatwhitespace=false,
    breaklines=true,
    captionpos=b, % Título siempre abajo
    keepspaces=true,
    numbers=left,                    
    numbersep=5pt,                  
    showspaces=false,                
    showstringspaces=false,
    showtabs=false,                  
    tabsize=2,
    frame=single,
    rulecolor=\color{grayborder},
    aboveskip=15pt,
    belowskip=10pt
}
\lstset{style=customcode}

% --- Definiciones de Lenguajes Adicionales no Nativos en Listings ---

% 1. Pseudocódigo (Español)
\lstdefinelanguage{Pseudocodigo}{
    sensitive=false,
    morekeywords={Inicio, Fin, Si, Entonces, Sino, FinSi, Para, Hasta, Hacer, FinPara, Mientras, FinMientras, Retornar, Funcion, Procedimiento, Escribir, Leer},
    morecomment=[l]{//},
    morestring=[b]"
}

% 1b. Pseudocode (English)
\lstdefinelanguage{Pseudocode}{
    sensitive=false,
    morekeywords={Begin, End, If, Then, Else, EndIf, For, To, Do, EndFor, While, EndWhile, Return, Function, Procedure, Write, Read},
    morecomment=[l]{//},
    morestring=[b]"
}

% 2. Rust
\lstdefinelanguage{Rust}{
    sensitive=true,
    morekeywords={fn, let, mut, impl, struct, enum, trait, use, mod, pub, match, return, if, else, loop, while, for, in, unsafe, async, await, type, as, self, Self},
    morecomment=[l]{//},
    morecomment=[s]{/*}{*/},
    morestring=[b]",
    morestring=[b]'
}

% 3. JavaScript
\lstdefinelanguage{JavaScript}{
    sensitive=true,
    morekeywords={const, let, var, function, return, if, else, for, while, class, extends, new, import, export, from, true, false, null, undefined, console, log},
    morecomment=[l]{//},
    morecomment=[s]{/*}{*/},
    morestring=[b]",
    morestring=[b]'
}

% 4. CSS
\lstdefinelanguage{CSS}{
    sensitive=false,
    morekeywords={color, background, margin, padding, border, display, flex, position, width, height, font, family, size, weight, align, justify},
    morecomment=[s]{/*}{*/},
    morestring=[b]"
}

% 5. Lua
\lstdefinelanguage{Lua}{
    sensitive=true,
    morekeywords={local, function, end, then, else, elseif, if, for, while, do, repeat, until, nil, true, false, return, break, in, pairs, ipairs},
    morecomment=[l]{--},
    morecomment=[s]{--[[}{]]},
    morestring=[b]",
    morestring=[b]'
}

% =========================================================================
% DATOS DEL DOCUMENTO
% =========================================================================
\title{\textbf{\color{primary}Plantilla General Académica con Ejemplos de Soporte}}
\author{Jair Aguilar Peña}
\date{\today}

% =========================================================================
% INICIO DEL DOCUMENTO
% =========================================================================
\begin{document}

\maketitle

\begin{abstract}
Este documento consolida la plantilla general académica para Ingeniería de Sistemas Cuantitativos. Muestra la integración de la tipografía Charter serif, colores corporativos y la configuración de leyendas independientes para figuras, tablas y códigos. Adicionalmente, incluye ejemplos listos de sintaxis para siete lenguajes de programación y pseudocódigo.
\end{abstract}

\tableofcontents
\newpage

\section{Estructura y Ejemplos de Leyendas (Captions)}

Las leyendas se formatean automáticamente con una jerarquía dedicada. Las etiquetas se resaltan en negrita y color verde azulado primario para asegurar un estándar editorial limpio.

\subsection{Ejemplo de Tabla de Alta Calidad}
Evite líneas verticales. Use el paquete \texttt{booktabs} como se muestra en la Tabla~\ref{tab:parametros}.

\begin{table}[h]
\centering
\caption{Parámetros de prueba del sistema en tiempo real.}
\label{tab:parametros}
\begin{tabular}{lcc}
    \toprule
    \textbf{Componente} & \textbf{Latencia Promedio} & \textbf{Desviación Estándar} \\
    \midrule
    Decodificador Binario & $120\text{ ns}$ & $\pm 5\text{ ns}$ \\
    Motor de Riesgo & $85\text{ ns}$ & $\pm 2\text{ ns}$ \\
    Enrutador de Órdenes & $150\text{ ns}$ & $\pm 8\text{ ns}$ \\
    \bottomrule
\end{tabular}
\end{table}

\subsection{Ejemplo de Figura o Imagen}
Las imágenes se integran usando el entorno estándar de figuras (ver Figura~\ref{fig:ejemplo}).

\begin{figure}[h]
\centering
% \includegraphics[width=0.4\textwidth]{nombre_imagen.png} % Reemplace con su imagen real
\framebox(200,80){Espacio reservado para Diagrama}
\caption{Esquema geométrico y flujo lógico del procesador.}
\label{fig:ejemplo}
\end{figure}

\newpage
\section{Ejemplos de Sintaxis de Programación}

\subsection{Pseudocódigo y Leyendas Dinámicas (Español / English)}

\begin{lstlisting}[language=Pseudocodigo, caption={Estructura básica de una búsqueda binaria en pseudocódigo.}]
Inicio
    Funcion BuscarElemento(arreglo, objetivo)
        izq = 0
        der = Longitud(arreglo) - 1
        Mientras izq <= der Hacer
            medio = (izq + der) / 2
            Si arreglo[medio] == objetivo Entonces
                Retornar medio
            Sino Si arreglo[medio] < objetivo Entonces
                izq = medio + 1
            Sino
                der = medio - 1
            FinSi
        FinMientras
        Retornar -1
    FinFuncion
Fin
\end{lstlisting}

\begin{otherlanguage}{english}
\begin{lstlisting}[language=Pseudocode, caption={Basic binary search structure in pseudocode.}]
Begin
    Function BinarySearch(array, target)
        left = 0
        right = Length(array) - 1
        While left <= right Do
            mid = (left + right) / 2
            If array[mid] == target Then
                Return mid
            Else If array[mid] < target Then
                left = mid + 1
            Else
                right = mid - 1
            EndIf
        EndWhile
        Return -1
    EndFunction
End
\end{lstlisting}
\end{otherlanguage}

\subsection{Python}
\begin{lstlisting}[language=Python, caption={Calculo de error de latencia en Python.}]
def calcular_error(medido: float, teorico: float) -> float:
    # Retorna el porcentaje de error absoluto
    assert teorico > 0, "El valor teórico debe ser mayor a cero"
    return abs(medido - teorico) / teorico * 100
\end{lstlisting}

\subsection{Rust}
\begin{lstlisting}[language=Rust, caption={Verificacion de propiedad y referencias en Rust.}]
fn main() {
    let mut data: String = String::from("Sistemas Cuantitativos");
    {
        let ref_immut = &data; // Préstamo inmutable
        println!("Lectura: {}", ref_immut);
    } // Termina el ámbito de la referencia inmutable
    
    let ref_mut = &mut data; // Préstamo mutable
    ref_mut.push_str(" (Ph.D. Track)");
    println!("Resultado: {}", ref_mut);
}
\end{lstlisting}

\subsection{C y C++}
\begin{lstlisting}[language=C, caption={Deteccion de endianness de la CPU en C.}]
#include <stdio.h>

int main(void) {
    unsigned int word = 0x01234567;
    unsigned char *byte_ptr = (unsigned char *)&word;
    
    if (*byte_ptr == 0x67) {
        printf("Arquitectura Little-Endian\n");
    } else {
        printf("Arquitectura Big-Endian\n");
    }
    return 0;
}
\end{lstlisting}

\subsection{LaTeX}
\begin{lstlisting}[language={[LaTeX]TeX}, caption={Fragmento de codigo matematico estructurado en LaTeX.}]
\begin{equation}
    A \vect{x} = \lambda \vect{x}
\end{equation}
Donde $\lambda$ representa el valor característico de la matriz de coeficientes.
\end{lstlisting}

\subsection{HTML, CSS y JavaScript (Tecnologías Web)}
\begin{lstlisting}[language=HTML, caption={Estructura básica e interactividad en HTML.}]
<div id="app">
    <h1>Quantitative Systems Engineering</h1>
    <button onclick="registrar()">Registrar Avance</button>
</div>
\end{lstlisting}

\begin{lstlisting}[language=CSS, caption={Estilo formal del bloque principal en CSS.}]
#app {
    background: #f7f7f5;
    border: 1px solid #d5d5d5;
    padding: 20px;
    margin: 15px auto;
}
\end{lstlisting}

\begin{lstlisting}[language=JavaScript, caption={Función de registro en JavaScript.}]
function registrar() {
    const timestamp = Date.now();
    console.log("Evento registrado a las: " + timestamp);
}
\end{lstlisting}

\subsection{Lua}
\begin{lstlisting}[language=Lua, caption={Script de configuracion para cargadores rápidos en Lua.}]
local config = {
    theme = "sobrio",
    latency_threshold = 200, -- nanosegundos
    verbose = true
}

function check_performance(value)
    if value > config.latency_threshold then
        return false
    else
        return true
    end
end
\end{lstlisting}

\end{document}
